% =============================================================================
% Introducción
% =============================================================================
\section*{Introducción}
\addcontentsline{toc}{section}{Introducción}

En la nueva era de la información digital en la que nos encontramos, el análisis de datos textuales es clave. Los motores de búsqueda evolucionan constantemente, migrando desde las búsquedas en base a coincidencia exacta de caracteres hacia sistemas capaces de entender el contexto, el significado y la relación conceptual entre los textos. 

\vspace{0.2cm}
Esta transición es la base para las herramientas y tecnologías actuales, como los asistentes virtuales, los sistemas de recomendación predictiva y las aplicaciones de inteligencia artificial.

\vspace{0.4cm}
Para que un sistema logre procesar e interpretar el lenguaje humano, es necesario modelar la información cualitativa en forma de estructuras matemáticas. El álgebra lineal proporciona las herramientas necesarias e ideales para este propósito. 

\vspace{0.2cm}
Al representar palabras, frases o documentos complejos como vectores dentro de espacios multidimensionales, podemos cuantificar los conceptos mediante propiedades geométricas abstractas. En este nuevo entorno, la similitud semántica entre dos elementos se traduce en la proximidad espacial o en la orientación angular de los respectivos vectores representativos.

\vspace{0.4cm}
El presente proyecto de investigación aborda el diseño, análisis e implementación computacional de un buscador semántico mediante la representación vectorial de texto. 

\vspace{0.2cm}
A través del modelo estructural Bag-of-Words (Bolsa de Palabras), se sistematiza la construcción de una matriz documento-término. Posteriormente, mediante la aplicación de operaciones algebraicas fundamentales —como el producto punto, el cálculo de la norma euclidiana y la similitud del coseno—, se desarrolla un motor de búsqueda capaz de procesar consultas en lenguaje natural y devolver los documentos más relevantes.

\vspace{0.4cm}
El documento contiene un análisis crítico de los resultados computacionales, explorando el impacto de la dimensión del espacio (tamaño del vocabulario), analizando la efectividad del sistema, y reflexionando sobre las ventajas y limitaciones de este enfoque con respecto a las implementaciones y herramientas actualmente utilizadas para el procesamiento de lenguaje natural.
